\documentclass[11pt]{article}
% ===========================================================================
%  Shared preamble for the AP Calculus AB exam-preparation worksheets.
%  Input right after \documentclass. Files beginning with "_" are not compiled
%  on their own. Navy and gold, to match the course's lesson pages.
% ===========================================================================
\usepackage[letterpaper,top=2.3cm,bottom=2.2cm,left=1.7cm,right=1.7cm,headheight=28pt,footskip=30pt]{geometry}
\usepackage{amsmath,amssymb}
\usepackage[T1]{fontenc}
\usepackage[utf8]{inputenc}
\usepackage{lmodern}
\usepackage{xcolor}
\usepackage[most]{tcolorbox}
\usepackage{enumitem}
\usepackage{fancyhdr}
\usepackage{array}
\usepackage{tikz}
\usepackage{pgfplots}
\pgfplotsset{compat=1.18}

\definecolor{navy}{HTML}{161E54}
\definecolor{ink}{HTML}{27357E}
\definecolor{gold}{HTML}{B98A00}
\definecolor{goldbg}{HTML}{FFFDF3}
\definecolor{inkbg}{HTML}{F6F7FD}
\definecolor{rule}{HTML}{C9CFF0}
\definecolor{okgreen}{HTML}{1B7A44}
\definecolor{okbg}{HTML}{F1F8F3}

\pagestyle{fancy}
\fancyhf{}
\fancyhead[L]{\footnotesize Name: \underline{\hspace{4cm}}}
\fancyhead[C]{\textbf{\color{navy}Integration Academy}\\[-2pt]{\scriptsize Calculus AB \textperiodcentered{} AP\textsuperscript{\textregistered} Exam Preparation}}
\fancyfoot[L]{\tiny\color{black!55}\parbox[t]{0.78\textwidth}{AP\textsuperscript{\textregistered} is a registered trademark of the College Board, which was not involved in the production of, and does not endorse, this material. \textcopyright{} Integration Academy.}}
\fancyfoot[R]{\scriptsize\thepage}
\renewcommand{\headrulewidth}{1.2pt}
\renewcommand{\headrule}{\hbox to\headwidth{\color{gold}\leaders\hrule height \headrulewidth\hfill}}
\renewcommand{\footrulewidth}{0pt}
\setlength{\parindent}{0pt}
\setlength{\parskip}{4pt}

% Title block: \wstitle{code}{title}{unit line}
\newcommand{\wstitle}[3]{\begin{center}{\small\bfseries\color{gold}WORKSHEET #1}\\[3pt]{\LARGE\bfseries\color{navy}#2}\\[3pt]{\small\color{black!60}#3}\end{center}\vspace{2pt}}
% Section heading
\newcommand{\sech}[1]{\par\vspace{10pt}{\large\bfseries\color{navy}#1}\par\nobreak\vspace{2pt}{\color{rule}\hrule height0.8pt}\vspace{7pt}}
% A level tag printed beside a group of questions
\newcommand{\level}[1]{\par\vspace{6pt}{\footnotesize\bfseries\color{gold}\MakeUppercase{#1}}\par\nobreak\vspace{2pt}}

\newtcolorbox{ideabox}{colframe=ink,colback=inkbg,coltitle=white,colbacktitle=ink,fonttitle=\bfseries\small,title={Key ideas},arc=2pt,boxrule=0.8pt,left=8pt,right=8pt}
\newtcolorbox{example}[1]{colframe=rule,colback=white,coltitle=navy,colbacktitle=inkbg,fonttitle=\bfseries\small,title={#1},arc=2pt,boxrule=0.8pt,borderline west={3pt}{0pt}{ink},left=10pt,right=8pt}
\newtcolorbox{frbox}[1]{colframe=gold,colback=goldbg,coltitle=white,colbacktitle=gold,fonttitle=\bfseries\small,title={#1},arc=2pt,boxrule=0.8pt,breakable,left=8pt,right=8pt}
\newtcolorbox{solbox}{colframe=okgreen,colback=okbg,coltitle=white,colbacktitle=okgreen,fonttitle=\bfseries\small,title={Worked solutions},arc=2pt,boxrule=0.8pt,breakable,left=8pt,right=8pt}
\newtcolorbox{scorebox}[1]{colframe=navy,colback=white,coltitle=white,colbacktitle=navy,fonttitle=\bfseries\small,title={#1},arc=2pt,boxrule=0.8pt,breakable,left=8pt,right=8pt}

\newcommand{\soln}{\par\smallskip\textbf{\color{ink}Solution.}\ }
\newcommand{\result}[1]{\par\smallskip\textbf{#1}}
% blank working space of a given height
\newcommand{\work}[1]{\par\vspace{#1}}
% a point value in the margin of a scoring guide
\newcommand{\pts}[1]{\hfill{\footnotesize\bfseries\color{gold}[#1]}}

% Coordinate grid for graph questions: \begin{apgraph}{xmin}{xmax}{ymin}{ymax} ... \end{apgraph}
\newenvironment{apgraph}[4]{\begin{center}\begin{tikzpicture}\begin{axis}[
  axis lines=middle, xmin=#1, xmax=#2, ymin=#3, ymax=#4,
  grid=both, grid style={black!12}, xtick distance=1, ytick distance=1,
  ticklabel style={font=\scriptsize}, xlabel={$x$}, ylabel={$y$},
  width=10cm, height=7.2cm, samples=80, clip=false,
  every axis plot/.append style={line width=1.1pt, color=ink}]}{\end{axis}\end{tikzpicture}\end{center}}
\newcommand{\closeddot}[2]{\addplot[only marks, mark=*, mark size=2.2pt, color=ink] coordinates {(#1,#2)};}
\newcommand{\opendot}[2]{\addplot[only marks, mark=*, mark size=2.2pt, color=ink, mark options={fill=white}] coordinates {(#1,#2)};}

% Multiple-choice options: four across, or stacked when the options are long.
\newcommand{\opts}[4]{\par\vspace{3pt}\begin{tabular}{@{}p{0.235\linewidth}p{0.235\linewidth}p{0.235\linewidth}p{0.235\linewidth}@{}}(A)\ #1 & (B)\ #2 & (C)\ #3 & (D)\ #4\end{tabular}\par\vspace{9pt}}
\newcommand{\optsv}[4]{\par\vspace{3pt}\begin{tabular}{@{}l@{\ \ }p{0.9\linewidth}@{}}(A) & #1\\(B) & #2\\(C) & #3\\(D) & #4\end{tabular}\par\vspace{9pt}}

\fancyhead[R]{\footnotesize\color{black!60}Worksheet 2.2}
\begin{document}

\wstitle{2.2}{The Derivative as a Limit}{Unit 2: The Derivative}

\begin{ideabox}
\begin{itemize}[leftmargin=*,itemsep=2pt,topsep=2pt]
  \item $f'(x)=\displaystyle\lim_{h\to 0}\dfrac{f(x+h)-f(x)}{h}$, wherever this limit exists.
  \item At a single point: $f'(a)=\displaystyle\lim_{x\to a}\dfrac{f(x)-f(a)}{x-a}$.
  \item Keep the symbol $\lim$ on every line until the limit has actually been taken.
  \item A limit that has the form of a difference quotient can be \textbf{recognised} as a derivative and evaluated with a rule.
\end{itemize}
\end{ideabox}

\sech{Worked examples}
\begin{example}{Example 1 \textperiodcentered{} A polynomial}
Use the definition to find $f'(x)$ for $f(x)=x^2-5x$.\soln $\displaystyle\lim_{h\to 0}\dfrac{(x+h)^2-5(x+h)-x^2+5x}{h}=\displaystyle\lim_{h\to 0}\dfrac{2xh+h^2-5h}{h}=\displaystyle\lim_{h\to 0}(2x+h-5)$.\result{$f'(x)=2x-5$.}
\end{example}

\begin{example}{Example 2 \textperiodcentered{} A square root}
Use the definition to find $f'(x)$ for $f(x)=\sqrt{2x+1}$.\soln Multiply by the conjugate: $\displaystyle\lim_{h\to 0}\dfrac{(2x+2h+1)-(2x+1)}{h\left(\sqrt{2x+2h+1}+\sqrt{2x+1}\right)}=\displaystyle\lim_{h\to 0}\dfrac{2}{\sqrt{2x+2h+1}+\sqrt{2x+1}}$.\result{$f'(x)=\dfrac{1}{\sqrt{2x+1}}$.}
\end{example}

\begin{example}{Example 3 \textperiodcentered{} Recognising a derivative}
Evaluate $\displaystyle\lim_{h\to 0}\dfrac{\sqrt{9+h}-3}{h}$.\soln This is $g'(9)$ for $g(x)=\sqrt{x}$, and $g'(x)=\dfrac{1}{2\sqrt x}$.\result{$\dfrac16$.}
\end{example}

\newpage
\sech{Practice}

\level{Core}
\begin{enumerate}[leftmargin=*,itemsep=2pt]
  \item Use the definition to find $f'(x)$ for $f(x)=3x+2$.\work{2.3cm}
  \item Use the definition to find $f'(x)$ for $f(x)=x^2+x$.\work{2.3cm}
  \item Evaluate $\displaystyle\lim_{h\to 0}\dfrac{(5+h)^2-25}{h}$.\work{2.3cm}
  \item Evaluate $\displaystyle\lim_{x\to 2}\dfrac{x^5-32}{x-2}$ by recognising it as a derivative.\work{2.3cm}
\end{enumerate}

\level{Exam level}
\begin{enumerate}[leftmargin=*,itemsep=2pt,start=5]
  \item Use the definition to find $f'(x)$ for $f(x)=\dfrac{1}{x^2}$.\work{2.8cm}
  \item Use the definition to find $f'(x)$ for $f(x)=x^3+2x$.\work{2.8cm}
  \item If $f'(4)=-3$, find $\displaystyle\lim_{h\to 0}\dfrac{f(4+h)-f(4)}{2h}$.\work{2.8cm}
  \item Evaluate $\displaystyle\lim_{h\to 0}\dfrac{\ln(2+h)-\ln2}{h}$.\work{2.8cm}
\end{enumerate}

\begin{frbox}{Question 9 \textperiodcentered{} Free-response practice}
Let $f(x)=\dfrac{3}{x-1}$ for $x\neq1$.
\begin{enumerate}[label=(\alph*),leftmargin=*,itemsep=2pt]
  \item Use the limit definition to find $f'(x)$.\work{2.6cm}
  \item Write an equation of the line tangent to the graph of $f$ at $x=2$.\work{2.6cm}
  \item Evaluate $\displaystyle\lim_{h\to 0}\dfrac{f(4+h)-f(4)}{h}$.\work{2.6cm}
\end{enumerate}
\end{frbox}

\newpage
\begin{solbox}
\begin{enumerate}[leftmargin=*,itemsep=4pt]
  \item $\displaystyle\lim_{h\to 0}\dfrac{3(x+h)+2-3x-2}{h}=\displaystyle\lim_{h\to 0}3=3$.
  \item $\displaystyle\lim_{h\to 0}\dfrac{2xh+h^2+h}{h}=2x+1$.
  \item The derivative of $x^2$ at 5: $10$.
  \item The derivative of $x^5$ at 2: $5(2)^4=80$.
  \item $\displaystyle\lim_{h\to 0}\dfrac{x^2-(x+h)^2}{h\,x^2(x+h)^2}=\displaystyle\lim_{h\to 0}\dfrac{-2x-h}{x^2(x+h)^2}=-\dfrac{2}{x^3}$.
  \item $\displaystyle\lim_{h\to 0}\dfrac{3x^2h+3xh^2+h^3+2h}{h}=3x^2+2$.
  \item $\dfrac12f'(4)=-\dfrac32$.
  \item The derivative of $\ln x$ at 2: $\dfrac12$.
  \item \begin{enumerate}[label=(\alph*),leftmargin=*,itemsep=3pt]
    \item $\displaystyle\lim_{h\to 0}\dfrac{\frac{3}{x+h-1}-\frac{3}{x-1}}{h}=\displaystyle\lim_{h\to 0}\dfrac{-3}{(x+h-1)(x-1)}=-\dfrac{3}{(x-1)^2}$.
    \item $f(2)=3$ and $f'(2)=-3$: $y-3=-3(x-2)$.
    \item This is $f'(4)=-\dfrac{3}{9}=-\dfrac13$.
  \end{enumerate}
\end{enumerate}
\end{solbox}

\end{document}
